\chapter{Proofs of the finite-sample estimates}
\label{app:finite-sample-proofs}

This appendix proves the finite-sample results from
Section~\ref{sec:results-finite-sample}. We work with the truncated Fourier subspace
\(\mathcal H_K\), the feature matrix \(\Phi\), the sampled kernel matrix \(A\), the Gram matrix
\(G\), and the matrix \(H\) defined in
\eqref{eq:method-feature-matrix}--\eqref{eq:method-compressed-operator}.

\section{Preliminaries}
\label{sec:app-preliminaries}

For each basis function in
\eqref{eq:method-fourier-basis-phi0}--\eqref{eq:method-fourier-basis-modes},
\begin{equation}
	|\phi_p(\varphi)|\le \sqrt2
	\qquad
	\text{for all }\varphi\in[0,2\pi),\quad 1\le p\le d.
	\label{eq:app-basis-bound}
\end{equation}
We also write
\begin{equation}
	\kappa:=\sup_{\theta\in[0,2\pi)}|\Theta(\theta)|<\infty.
	\label{eq:app-kappa}
\end{equation}

For each basis index \(p\), let \(\lambda_p\) denote the eigenvalue of the operator \(T\)
associated with \(\phi_p\), so that
\[
	T\phi_p=\lambda_p \phi_p.
\]
Define the finite-sample corrected eigenvalues by
\begin{equation}
	\lambda_p^{(n)}
	=
	\left(1-\frac{1}{n}\right)\lambda_p+\frac{\Theta(0)}{n},
	\label{eq:app-finite-sample-corrected-eigs}
\end{equation}
and let
\[
	\Lambda^{(n)}=\operatorname{diag}(\lambda_p^{(n)}).
\]

\section{Expectation of \(H\)}
\label{sec:app-expected-H}

\begin{proof}[Proof of Proposition~\ref{prop:results-expectation-H}]
	By definition,
	\begin{equation}
		H_{pq}
		=
		\frac{1}{n^2}\sum_{i=1}^n\sum_{j=1}^n
		\phi_p(\varphi_i)\Theta(\varphi_i-\varphi_j)\phi_q(\varphi_j).
		\label{eq:app-H-entry}
	\end{equation}
	Split the sum into the cases \(i\neq j\) and \(i=j\):
	\[
		H_{pq}
		=
		\frac{1}{n^2}\sum_{i\neq j}
		\phi_p(\varphi_i)\Theta(\varphi_i-\varphi_j)\phi_q(\varphi_j)
		+
		\frac{1}{n^2}\sum_{i=1}^n
		\phi_p(\varphi_i)\Theta(0)\phi_q(\varphi_i).
	\]

	For the diagonal terms,
	\[
		\mathbb E\!\left[
			\frac{1}{n^2}\sum_{i=1}^n
			\phi_p(\varphi_i)\Theta(0)\phi_q(\varphi_i)
			\right]
		=
		\frac{\Theta(0)}{n}\delta_{pq}.
	\]

	For the off-diagonal terms, independence of \(\varphi_i\) and \(\varphi_j\) gives
	\begin{align*}
		\mathbb E\!\left[
			\phi_p(\varphi_i)\Theta(\varphi_i-\varphi_j)\phi_q(\varphi_j)
			\right]
		 & =
		\int_0^{2\pi}\phi_p(\varphi)
		\left(
		\int_0^{2\pi}\Theta(\varphi-\psi)\phi_q(\psi)\,d\mu(\psi)
		\right)
		d\mu(\varphi) \\
		 & =
		\int_0^{2\pi}\phi_p(\varphi)(T\phi_q)(\varphi)\,d\mu(\varphi)
		=
		\lambda_q\delta_{pq}.
	\end{align*}
	Since there are \(n(n-1)\) off-diagonal pairs, we conclude that
	\[
		\mathbb E[H_{pq}]
		=
		\left(1-\frac{1}{n}\right)\lambda_q\delta_{pq}
		+
		\frac{\Theta(0)}{n}\delta_{pq}
		=
		\lambda_q^{(n)}\delta_{pq}.
	\]
	Therefore \(\mathbb E[H]=\Lambda^{(n)}\).
\end{proof}

\section{Concentration of the Gram matrix}
\label{sec:app-gram-proof}

\begin{proof}[Proof of \eqref{eq:results-gram-concentration}]
	For fixed \(p,q\),
	\[
		G_{pq}
		=
		\frac{1}{n}\sum_{i=1}^n Z_i,
		\qquad
		Z_i:=\phi_p(\varphi_i)\phi_q(\varphi_i).
	\]
	The variables \(Z_i\) are independent and identically distributed, with
	\[
		\mathbb E[Z_i]=\delta_{pq}
	\]
	by orthonormality of the Fourier basis.

	Set
	\[
		Y_i:=Z_i-\delta_{pq}.
	\]
	Then \(\mathbb E[Y_i]=0\) and
	\[
		G_{pq}-\delta_{pq}=\frac{1}{n}\sum_{i=1}^n Y_i.
	\]

	By \eqref{eq:app-basis-bound},
	\[
		|Z_i|\le 2.
	\]
	If \(p\neq q\), then \(\delta_{pq}=0\), so \(|Y_i|=|Z_i|\le 2\). If \(p=q\), then
	\[
		Y_i=\phi_p(\varphi_i)^2-1,
	\]
	and since \(0\le \phi_p(\varphi_i)^2\le 2\), we have \(|Y_i|\le 1\le 2\). Thus
	\[
		|Y_i|\le 2
	\]
	for all \(i\).

	Next,
	\[
		\mathrm{Var}(Y_i)\le \mathbb E[Y_i^2]\le \mathbb E[Z_i^2].
	\]
	Using again \eqref{eq:app-basis-bound},
	\[
		Z_i^2=\phi_p(\varphi_i)^2\phi_q(\varphi_i)^2\le 2\,\phi_p(\varphi_i)^2.
	\]
	Taking expectations and using orthonormality gives
	\[
		\mathbb E[Z_i^2]\le 2\int \phi_p(\varphi)^2\,d\mu(\varphi)=2.
	\]
	Hence
	\[
		\mathrm{Var}(Y_i)\le 2.
	\]

	We now apply Bernstein's inequality \citep{Vershynin_2018a} in the following form: if \(Y_1,\dots,Y_n\) are independent,
	centered random variables such that
	\[
		|Y_i|\le M,
		\qquad
		\mathrm{Var}(Y_i)\le \sigma^2
	\]
	for every \(i\), then for every \(\varepsilon>0\),
	\[
		\mathbb P\!\left(
		\left|\frac{1}{n}\sum_{i=1}^n Y_i\right|
		\ge \varepsilon
		\right)
		\le
		2\exp\!\left(
		-\frac{n\varepsilon^2}{2\sigma^2+\frac{2}{3}M\varepsilon}
		\right).
	\]

	In our case, we have already shown that \(\mathbb E[Y_i]=0\), \(|Y_i|\le 2\), and
	\(\mathrm{Var}(Y_i)\le 2\). Thus we may take
	\[
		M=2,
		\qquad
		\sigma^2=2.
	\]
	Since
	\[
		G_{pq}-\delta_{pq}=\frac{1}{n}\sum_{i=1}^n Y_i,
	\]

	Bernstein's inequality for the sample mean yields
	\[
		\mathbb P\!\left(|G_{pq}-\delta_{pq}|\ge \varepsilon\right)
		\le
		2\exp\!\left(
		-\frac{n\varepsilon^2}{4+\frac43\varepsilon}
		\right).
	\]

	Since \(G-I_d\) is symmetric, it suffices to union bound over the \(d(d+1)/2\) entries in
	the upper triangle. Therefore
	\[
		\mathbb P\!\left(\|G-I_d\|_{\max}\ge \varepsilon\right)
		\le
		d(d+1)\exp\!\left(
		-\frac{n\varepsilon^2}{4+\frac43\varepsilon}
		\right),
	\]
	which proves \eqref{eq:results-gram-concentration}.
\end{proof}

\section{Concentration of \(H\)}
\label{sec:app-H-proof}

\begin{proof}[Proof of \eqref{eq:results-H-concentration}]
	Fix \(p,q\), and write \(H_{pq}=F(\varphi_1,\dots,\varphi_n)\), where
	\[
		F(\varphi_1,\dots,\varphi_n)
		=
		\frac{1}{n^2}\sum_{i=1}^n\sum_{j=1}^n
		\phi_p(\varphi_i)\Theta(\varphi_i-\varphi_j)\phi_q(\varphi_j).
	\]

	We apply McDiarmid's inequality \citep{Vershynin_2018b} in the following form: if
	\(X_1,\dots,X_n\) are independent and
	\(F(X_1,\dots,X_n)\) satisfies the bounded-differences condition
	\[
		|F(x_1,\dots,x_s,\dots,x_n)-F(x_1,\dots,x_s',\dots,x_n)|\le c_s
	\]
	for each \(s\), then for every \(\varepsilon>0\),
	\[
		\mathbb P\!\left(|F-\mathbb EF|\ge \varepsilon\right)
		\le
		2\exp\!\left(-\frac{2\varepsilon^2}{\sum_{s=1}^n c_s^2}\right).
	\]

	It therefore remains to estimate the effect of changing a single sample point, say
	\(\varphi_s\mapsto \varphi_s'\), while keeping all others fixed.
	Only those summands in the double sum with \(i=s\) or \(j=s\) can change. There are at most
	\(2n\) such summands.

	By \eqref{eq:app-basis-bound} and \eqref{eq:app-kappa}, each summand satisfies
	\[
		|\phi_p(\cdot)\Theta(\cdot)\phi_q(\cdot)|
		\le
		\sqrt2\,\kappa\,\sqrt2
		=
		2\kappa.
	\]
	Hence, when \(\varphi_s\) is replaced by \(\varphi_s'\), a single summand can change by at most \(4\kappa\).

	Because at most \(2n\) summands are affected, the total change in the double sum is bounded by
	\[
		(2n)\cdot (4\kappa)=8\kappa n.
	\]
	Finally, since \(F\) contains the prefactor \(1/n^2\), the change in \(F\) itself is bounded by
	\[
		\frac{8\kappa n}{n^2}=\frac{8\kappa}{n}.
	\]
	Thus we may take
	\[
		c_s=\frac{8\kappa}{n}
		\qquad\text{for all } s=1,\dots,n.
	\]

	Therefore
	\[
		\sum_{s=1}^n c_s^2
		=
		n\left(\frac{8\kappa}{n}\right)^2
		=
		\frac{64\kappa^2}{n}.
	\]
	McDiarmid's inequality yields
	\[
		\mathbb P\!\left(|H_{pq}-\mathbb E[H_{pq}]|\ge \varepsilon\right)
		\le
		2\exp\!\left(
		-\frac{2\varepsilon^2}{64\kappa^2/n}
		\right)
		=
		2\exp\!\left(-\frac{n\varepsilon^2}{32\kappa^2}\right).
	\]

	A union bound over all \(d^2\) entries gives
	\[
		\mathbb P\!\left(\|H-\Lambda^{(n)}\|_{\max}\ge \varepsilon\right)
		\le
		2d^2\exp\!\left(-\frac{n\varepsilon^2}{32\kappa^2}\right),
	\]
	as claimed.
\end{proof}
\begin{remark}
	The entry \(H_{pq}\) is a bounded second-order V-statistic, so Bernstein-type refinements are
	in principle available through concentration inequalities for bounded U/V-statistics. Since the
	McDiarmid bound already captures the correct \(n^{-1/2}\) concentration scale and is sufficient
	for the results in this thesis, we do not pursue this refinement here.
\end{remark}

\section{Control of the difference \(A\Phi-\Phi\Lambda^{(n)}\)}
\label{sec:app-APhi-proof}

\begin{proof}[Proof of \eqref{eq:results-APhi-concentration}]
	Fix \(i,p\). By definition,
	\[
		(A\Phi)_{i,p}
		=
		\frac{1}{n}\sum_{j=1}^n \Theta(\varphi_i-\varphi_j)\phi_p(\varphi_j).
	\]
	Split off the diagonal term \(j=i\):
	\[
		(A\Phi)_{i,p}
		=
		\frac{\Theta(0)}{n}\phi_p(\varphi_i)
		+
		\frac{1}{n}\sum_{j\neq i}\Theta(\varphi_i-\varphi_j)\phi_p(\varphi_j).
	\]

	Condition on \(\varphi_i\). For \(j\neq i\), define
	\[
		X_j:=\Theta(\varphi_i-\varphi_j)\phi_p(\varphi_j).
	\]
	Then \(\{X_j\}_{j\neq i}\) are conditionally independent and identically distributed.

	Their conditional mean is
	\[
		\mathbb E[X_j\mid \varphi_i]
		=
		\int_0^{2\pi}\Theta(\varphi_i-\psi)\phi_p(\psi)\,d\mu(\psi)
		=
		(T\phi_p)(\varphi_i)
		=
		\lambda_p\phi_p(\varphi_i).
	\]
	Therefore
	\[
		\mathbb E[(A\Phi)_{i,p}\mid \varphi_i]
		=
		\frac{\Theta(0)}{n}\phi_p(\varphi_i)
		+
		\frac{n-1}{n}\lambda_p\phi_p(\varphi_i)
		=
		\lambda_p^{(n)}\phi_p(\varphi_i).
	\]

	Now set
	\[
		Y_j:=X_j-\mathbb E[X_j\mid \varphi_i].
	\]
	Conditional on \(\varphi_i\), the variables \(\{Y_j\}_{j\neq i}\) are independent,
	centered, and identically distributed.

	We apply Bernstein's inequality \citep{Vershynin_2018a} in the following form: if \(Y_1,\dots,Y_m\) are independent,
	centered random variables such that
	\[
		|Y_j|\le M,
		\qquad
		\mathrm{Var}(Y_j)\le \sigma^2
	\]
	for all \(j\), then for every \(\eta>0\),
	\[
		\mathbb P\!\left(
		\left|
		\frac1m\sum_{j=1}^m Y_j
		\right|
		\ge \eta
		\right)
		\le
		2\exp\!\left(
		-\frac{m\eta^2}{2\sigma^2+\frac23 M\eta}
		\right).
	\]

	We now verify the assumptions conditionally on \(\varphi_i\).

	First, by \eqref{eq:app-basis-bound} and \eqref{eq:app-kappa},
	\[
		|X_j|
		\le
		\kappa \sqrt2
		\le
		2\kappa.
	\]
	Also,
	\[
		|\mathbb E[X_j\mid \varphi_i]|
		=
		\left|
		\int_0^{2\pi}\Theta(\varphi_i-\psi)\phi_p(\psi)\,d\mu(\psi)
		\right|
		\le
		\kappa \int_0^{2\pi}|\phi_p(\psi)|\,d\mu(\psi)
		\le
		\kappa \|\phi_p\|_{L^2(\mu)}
		=
		\kappa.
	\]
	Hence
	\[
		|Y_j|
		=
		|X_j-\mathbb E[X_j\mid \varphi_i]|
		\le
		2\kappa+\kappa
		=
		3\kappa.
	\]

	Next,
	\[
		\mathrm{Var}(Y_j\mid \varphi_i)
		=
		\mathrm{Var}(X_j\mid \varphi_i)
		\le
		\mathbb E[X_j^2\mid \varphi_i].
	\]
	Using again \eqref{eq:app-basis-bound} and \eqref{eq:app-kappa},
	\[
		\mathbb E[X_j^2\mid \varphi_i]
		=
		\int_0^{2\pi}\Theta(\varphi_i-\psi)^2\phi_p(\psi)^2\,d\mu(\psi)
		\le
		\kappa^2 \int_0^{2\pi}\phi_p(\psi)^2\,d\mu(\psi)
		=
		\kappa^2.
	\]
	Thus, conditionally on \(\varphi_i\), Bernstein's inequality applies with
	\[
		M=3\kappa,
		\qquad
		\sigma^2=\kappa^2.
	\]

	Let \(m:=n-1\). Applying Bernstein to the conditional sample mean
	\[
		\frac1m\sum_{j\neq i}Y_j
		=
		\frac1m\sum_{j\neq i}X_j-\lambda_p\phi_p(\varphi_i),
	\]
	we obtain
	\[
		\mathbb P\!\left(
		\left|
		\frac1m\sum_{j\neq i}X_j-\lambda_p\phi_p(\varphi_i)
		\right|
		\ge \eta
		\;\middle|\;
		\varphi_i
		\right)
		\le
		2\exp\!\left(
		-\frac{m\eta^2}{2\kappa^2+2\kappa\eta}
		\right).
	\]

	Now
	\[
		(A\Phi)_{i,p}-\lambda_p^{(n)}\phi_p(\varphi_i)
		=
		\frac{m}{n}
		\left(
		\frac1m\sum_{j\neq i}X_j-\lambda_p\phi_p(\varphi_i)
		\right).
	\]
	Thus the event
	\[
		\left|
		(A\Phi)_{i,p}-\lambda_p^{(n)}\phi_p(\varphi_i)
		\right|
		\ge \varepsilon
	\]
	implies
	\[
		\left|
		\frac1m\sum_{j\neq i}X_j-\lambda_p\phi_p(\varphi_i)
		\right|
		\ge \frac{n}{m}\varepsilon.
	\]
	Hence
	\[
		\mathbb P\!\left(
		\left|
		(A\Phi)_{i,p}-\lambda_p^{(n)}\phi_p(\varphi_i)
		\right|
		\ge \varepsilon
		\;\middle|\;
		\varphi_i
		\right)
		\le
		2\exp\!\left(
		-\frac{m\left(\frac{n}{m}\varepsilon\right)^2}
		{2\kappa^2+2\kappa\left(\frac{n}{m}\varepsilon\right)}
		\right).
	\]
	Removing the conditioning and simplifying gives
	\[
		\mathbb P\!\left(
		\left|
		(A\Phi)_{i,p}-\lambda_p^{(n)}\phi_p(\varphi_i)
		\right|
		\ge \varepsilon
		\right)
		\le
		2\exp\!\left(
		-\frac{n^2\varepsilon^2}{2\kappa^2(n-1)+2\kappa n\varepsilon}
		\right).
	\]
	Since \(n-1\le n\), we obtain
	\[
		\mathbb P\!\left(
		\left|
		(A\Phi)_{i,p}-\lambda_p^{(n)}\phi_p(\varphi_i)
		\right|
		\ge \varepsilon
		\right)
		\le
		2\exp\!\left(
		-\frac{n\varepsilon^2}{2\kappa^2+2\kappa\varepsilon}
		\right).
	\]

	Finally, taking a union bound over all \(nd\) pairs \((i,p)\) yields
	\[
		\mathbb P\!\left(\|A\Phi-\Phi\Lambda^{(n)}\|_{\max}\ge \varepsilon\right)
		\le
		2nd\exp\!\left(-\frac{n\varepsilon^2}{2\kappa^2+2\kappa\varepsilon}\right),
	\]
	which proves \eqref{eq:results-APhi-concentration}.
\end{proof}

\section{Proof of the simultaneous high-probability bound}
\label{sec:app-simultaneous-proof}

\begin{corollary}[High-probability bounds]
	\label{cor:results-simultaneous-control}
	Let \(\delta\in(0,1)\), and define
	\begin{equation}
		u_{n,d,\delta}^{(G)}
		:=
		\log\!\left(\frac{3d(d+1)}{\delta}\right),
		\label{eq:results-uG}
	\end{equation}
	\begin{equation}
		u_{n,d,\delta}^{(H)}
		:=
		\log\!\left(\frac{6d^2}{\delta}\right),
		\label{eq:results-uH}
	\end{equation}
	and
	\begin{equation}
		u_{n,d,\delta}^{(A)}
		:=
		\log\!\left(\frac{6nd}{\delta}\right).
		\label{eq:results-uA}
	\end{equation}
	Now set
	\begin{equation}
		\alpha_{n,d,\delta}
		:=
		2\sqrt{\frac{u_{n,d,\delta}^{(G)}}{n}}
		+\frac{2}{3}\frac{u_{n,d,\delta}^{(G)}}{n},
		\label{eq:results-alpha}
	\end{equation}
	\begin{equation}
		\beta_{n,d,\delta}
		:=
		\kappa\sqrt{\frac{32}{n}u_{n,d,\delta}^{(H)}},
		\label{eq:results-beta}
	\end{equation}
	and
	\begin{equation}
		\gamma_{n,d,\delta}
		:=
		\kappa\sqrt{\frac{2}{n}u_{n,d,\delta}^{(A)}}
		+\kappa\frac{u_{n,d,\delta}^{(A)}}{n}.
		\label{eq:results-gamma}
	\end{equation}
	Then, with probability at least \(1-\delta\),
	\begin{equation}
		\|G-I_d\|_{\max}\le \alpha_{n,d,\delta},
		\qquad
		\|H-\Lambda^{(n)}\|_{\max}\le \beta_{n,d,\delta},
		\qquad
		\|A\Phi-\Phi\Lambda^{(n)}\|_{\max}\le \gamma_{n,d,\delta}.
		\label{eq:results-simultaneous-bounds}
	\end{equation}
\end{corollary}

\begin{proof}[Proof of Corollary~\ref{cor:results-simultaneous-control}]
	Set
	\[
		u_{n,d,\delta}^{(G)}
		=
		\log\!\left(\frac{3d(d+1)}{\delta}\right),
		\qquad
		u_{n,d,\delta}^{(H)}
		=
		\log\!\left(\frac{6d^2}{\delta}\right),
		\qquad
		u_{n,d,\delta}^{(A)}
		=
		\log\!\left(\frac{6nd}{\delta}\right).
	\]
	Define
	\[
		\alpha_{n,d,\delta}
		=
		2\sqrt{\frac{u_{n,d,\delta}^{(G)}}{n}}
		+
		\frac{2}{3}\frac{u_{n,d,\delta}^{(G)}}{n},
	\]
	\[
		\beta_{n,d,\delta}
		=
		\kappa\sqrt{\frac{32}{n}u_{n,d,\delta}^{(H)}},
		\qquad
		\gamma_{n,d,\delta}
		=
		\kappa\sqrt{\frac{2}{n}u_{n,d,\delta}^{(A)}}
		+
		\kappa\frac{u_{n,d,\delta}^{(A)}}{n}.
	\]

	We first recall the standard Bernstein inversion: if
	\[
		\mathbb P\!\left(\left|\frac{1}{n}\sum_{i=1}^n Y_i\right|\ge t\right)
		\le
		2\exp\!\left(
		-\frac{nt^2}{2\sigma^2+\frac{2}{3}Mt}
		\right),
	\]
	then choosing
	\[
		t=\sqrt{\frac{2\sigma^2 u}{n}}+\frac{Mu}{3n}
	\]
	ensures
	\[
		\mathbb P\!\left(\left|\frac{1}{n}\sum_{i=1}^n Y_i\right|\ge t\right)\le 2e^{-u}.
	\]

	For the Gram matrix bound \eqref{eq:results-gram-concentration}, the Bernstein proof used
	\[
		\sigma^2=2,
		\qquad
		M=2.
	\]
	Thus, with
	\[
		u=u_{n,d,\delta}^{(G)},
	\]
	the choice
	\[
		\alpha_{n,d,\delta}
		=
		\sqrt{\frac{2\cdot 2\,u}{n}}+\frac{2u}{3n}
		=
		2\sqrt{\frac{u}{n}}+\frac{2u}{3n}
	\]
	yields
	\[
		\mathbb P\!\left(\|G-I_d\|_{\max}\ge \alpha_{n,d,\delta}\right)
		\le
		d(d+1)e^{-u_{n,d,\delta}^{(G)}}
		=
		\frac{\delta}{3}.
	\]

	For the compressed operator, the McDiarmid estimate \eqref{eq:results-H-concentration} gives
	\[
		\mathbb P\!\left(\|H-\Lambda^{(n)}\|_{\max}\ge \varepsilon\right)
		\le
		2d^2\exp\!\left(-\frac{n\varepsilon^2}{32\kappa^2}\right).
	\]
	Substituting \(\varepsilon=\beta_{n,d,\delta}\) gives
	\[
		\frac{n\beta_{n,d,\delta}^2}{32\kappa^2}
		=
		u_{n,d,\delta}^{(H)},
	\]
	and therefore
	\[
		\mathbb P\!\left(\|H-\Lambda^{(n)}\|_{\max}\ge \beta_{n,d,\delta}\right)
		\le
		2d^2e^{-u_{n,d,\delta}^{(H)}}
		=
		\frac{\delta}{3}.
	\]

	For the approximate eigen-action bound \eqref{eq:results-APhi-concentration}, the Bernstein proof
	used
	\[
		\sigma^2=\kappa^2,
		\qquad
		M=3\kappa.
	\]
	Hence, with
	\[
		u=u_{n,d,\delta}^{(A)},
	\]
	the standard Bernstein inversion gives the threshold
	\[
		\sqrt{\frac{2\kappa^2 u}{n}}+\frac{3\kappa u}{3n}
		=
		\kappa\sqrt{\frac{2u}{n}}+\kappa\frac{u}{n},
	\]
	which is precisely \(\gamma_{n,d,\delta}\). Therefore
	\[
		\mathbb P\!\left(\|A\Phi-\Phi\Lambda^{(n)}\|_{\max}\ge \gamma_{n,d,\delta}\right)
		\le
		2nd\,e^{-u_{n,d,\delta}^{(A)}}
		=
		\frac{\delta}{3}.
	\]

	A final union bound gives
	\[
		\mathbb P\!\left(
		\|G-I_d\|_{\max}\le \alpha_{n,d,\delta},\;
		\|H-\Lambda^{(n)}\|_{\max}\le \beta_{n,d,\delta},\;
		\|A\Phi-\Phi\Lambda^{(n)}\|_{\max}\le \gamma_{n,d,\delta}
		\right)
		\ge
		1-\delta,
	\]
	which is exactly \eqref{eq:results-simultaneous-bounds}.
\end{proof}

% \section{Proof of the corollary on the restricted sampled operator}
% \label{sec:app-corollary-proof}

% \begin{proof}[Proof of Corollary~\ref{cor:results-restricted-operator}]
% 	Recall that
% 	\[
% 		C:=G^{-1}H.
% 	\]
% 	We work on the event of Corollary~\ref{cor:results-simultaneous-control}, namely
% 	\[
% 		\|G-I_d\|_{\max}\le \alpha_{n,d,\delta},
% 		\qquad
% 		\|H-\Lambda^{(n)}\|_{\max}\le \beta_{n,d,\delta}.
% 	\]
%
% 	Passing from the entrywise max norm to the induced \(\ell^\infty\)-operator norm, we use
% 	\[
% 		\|M\|_\infty\le d\,\|M\|_{\max}
% 	\]
% 	for every \(d\times d\) matrix \(M\). Hence
% 	\begin{equation}
% 		\|G-I_d\|_\infty\le d\,\alpha_{n,d,\delta},
% 		\qquad
% 		\|H-\Lambda^{(n)}\|_\infty\le d\,\beta_{n,d,\delta}.
% 		\label{eq:app-norm-conversion}
% 	\end{equation}
%
% 	If \(d\,\alpha_{n,d,\delta}<1\), then \(G\) is invertible. Indeed,
% 	\[
% 		G = I_d - (I_d-G),
% 	\]
% 	and since \(\|I_d-G\|_\infty<1\), the Neumann-series argument gives
% 	\begin{equation}
% 		\|G^{-1}\|_\infty
% 		\le
% 		\frac{1}{1-\|I_d-G\|_\infty}
% 		\le
% 		\frac{1}{1-d\,\alpha_{n,d,\delta}}.
% 		\label{eq:app-G-inverse-bound}
% 	\end{equation}
%
% 	Using
% 	\begin{equation}
% 		C-\Lambda^{(n)}
% 		=
% 		G^{-1}H-\Lambda^{(n)}
% 		=
% 		G^{-1}\bigl(H-G\Lambda^{(n)}\bigr)
% 		=
% 		G^{-1}\bigl((H-\Lambda^{(n)})-(G-I_d)\Lambda^{(n)}\bigr),
% 		\label{eq:app-C-minus-Lambda}
% 	\end{equation}
% 	we obtain
% 	\[
% 		\|C-\Lambda^{(n)}\|_\infty
% 		\le
% 		\|G^{-1}\|_\infty
% 		\left(
% 		\|H-\Lambda^{(n)}\|_\infty
% 		+
% 		\|G-I_d\|_\infty\,\|\Lambda^{(n)}\|_\infty
% 		\right).
% 	\]
%
% 	It remains to bound \(\|\Lambda^{(n)}\|_\infty\). Since \(\Lambda^{(n)}\) is diagonal,
% 	\[
% 		\|\Lambda^{(n)}\|_\infty
% 		=
% 		\max_p |\lambda_p^{(n)}|.
% 	\]
% 	Now
% 	\[
% 		\lambda_p^{(n)}
% 		=
% 		\left(1-\frac1n\right)\lambda_p+\frac1n\Theta(0),
% 	\]
% 	and both \(|\lambda_p|\le \kappa\) and \(|\Theta(0)|\le \kappa\). Therefore
% 	\[
% 		|\lambda_p^{(n)}|
% 		\le
% 		\left(1-\frac1n\right)\kappa+\frac1n\kappa
% 		=
% 		\kappa,
% 	\]
% 	so
% 	\[
% 		\|\Lambda^{(n)}\|_\infty\le \kappa.
% 	\]
%
% 	Substituting \eqref{eq:app-norm-conversion} and \eqref{eq:app-G-inverse-bound} gives
% 	\[
% 		\|C-\Lambda^{(n)}\|_\infty
% 		\le
% 		\frac{d\,\beta_{n,d,\delta}+d\,\kappa\,\alpha_{n,d,\delta}}
% 		{1-d\,\alpha_{n,d,\delta}},
% 	\]
% 	which is exactly \eqref{eq:results-restricted-operator-bound}.
% \end{proof}
